% =========================================================
% Open-source Statement
% =========================================================
% This is an unofficial LaTeX template for the Research Proposal Report
% and Dissertation Implementation Plan of professional degree graduate
% students at Southeast University.
%
% This template is for academic typesetting and personal use only. It is
% not an official template released by Southeast University. Please check
% the latest official requirements before formal submission.
%
% License: MIT
% Author: Zhaodong Lyu
% =========================================================


\documentclass[UTF8,a4paper,zihao=-4]{ctexart}

\usepackage{fontspec}
\setmainfont{Times New Roman}
\setmonofont{Times New Roman}

\usepackage{geometry}
\usepackage{array}
\usepackage{tabularx}
\usepackage{longtable}
\usepackage{makecell}
\usepackage{multirow}
\usepackage{enumitem}
\usepackage{fancyhdr}
\usepackage{titlesec}
\usepackage[most]{tcolorbox}
\usepackage{lipsum}
\usepackage[numbers,sort&compress]{natbib}
\renewcommand{\bibsection}{}


\geometry{left=2.0cm,right=2.0cm,top=2.0cm,bottom=2.0cm}
\pagestyle{fancy}
\fancyhf{}
\cfoot{\thepage}
\renewcommand{\headrulewidth}{0pt}
\setlength{\parindent}{2em}
\setlength{\parskip}{0pt}
\renewcommand{\arraystretch}{1.35}
\titleformat{\section}{\centering\bfseries\zihao{4}}{\thesection}{0.5em}{}

\newcommand{\ulinebox}[2][4.8cm]{\underline{\makebox[#1]{#2}}}
\newcommand{\emptyline}[1][4.5cm]{\underline{\makebox[#1]{}}}
\newcommand{\coverfield}[2]{\makebox[5.6em][s]{#1}\quad \ulinebox[5.0cm]{#2}\\[1.55em]}

\newcolumntype{Y}{>{\centering\arraybackslash}X}
\newcolumntype{L}{>{\raggedright\arraybackslash}X}
\newcolumntype{C}[1]{>{\centering\arraybackslash}m{#1}}
\newcolumntype{P}[1]{>{\centering\arraybackslash}p{#1}}

\newcommand{\vlabel}[2]{%
\begin{minipage}[c][#1][c]{\linewidth}
\centering #2
\end{minipage}%
}

\tcbset{
    reportbox/.style={
        enhanced,
        breakable,
        width=\textwidth,
        colback=white,
        colframe=black,
        boxrule=0.4pt,
        arc=0pt,
        outer arc=0pt,
        boxsep=0pt,
        left=0pt,
        right=0pt,
        top=0pt,
        bottom=0pt,
        before skip=-\arrayrulewidth,
        after skip=0pt
    }
}

\begin{document}

\thispagestyle{empty}
\vspace*{0.1cm}
\begin{flushright}
{\zihao{-4} 学号：\ulinebox[2.2cm]{xxxxxx}}\hspace*{0.5cm}
\end{flushright}

\vspace{1.7cm}
\begin{center}
{\zihao{2}\bfseries 东南大学专业学位型研究生}\\[0.9em]
{\zihao{2}\bfseries 学位论文开题报告及论文实施计划}
\end{center}

\vspace{2.25cm}
\begin{flushright}
    \centering
    \begin{minipage}{9.0cm}
    \zihao{4}
    \coverfield{院（系、所）}{xxxx学院}
    \coverfield{学位类别}{工程博士}
    \coverfield{专业领域}{xxxx}
    \coverfield{研究生姓名}{xxx}
    \coverfield{指导教师（校内）}{xxx}
    \coverfield{指导教师（校外）}{xxx}
    \coverfield{开题报告日期}{202x 年 xx 月}
    \end{minipage}
    \hspace*{1.15cm}
\end{flushright}

\vfill
\begin{center}
{\zihao{4} 东南大学研究生院制表}
\end{center}

\newpage
\thispagestyle{empty}
\vspace*{0.15cm}
\begin{center}
{\zihao{3}\bfseries 填表须知}
\end{center}
\vspace{1cm}

\begin{enumerate}[leftmargin=1.8em,label=\arabic*、,itemsep=0.45em]
    \item 论文开题报告由研究生本人向审议小组报告并听取意见后，由研究生本人填写此表。
    \item 论文开题报告填写完成后，必须经导师审批，通过后方能提交。
    \item 博士生应在第四学期内、硕士生应在第三学期内完成此开题报告。开题报告经研究生秘书在网上审核确认（硕士生至少半年、博士生至少一年）后方可申请答辩。
    \item 研究生开题前应填写查新报告。查新报告对专业学位博士作为必要环节。博士生查新工作可委托图书馆负责，也可在完成网络文献检索类研究生课程的学习或参加学校组织的网络文献检索培训后，自行组织查新检索，自行组织查新需要详细文献查新述评作为附件。自行查新报告须经导师审查后由开题报告审核专家组审核签字（或盖章）。硕士生开题查新参考上述办法，不作硬性要求。
    \item 本表一式两份，一份研究生自留放入本人“研究生档案材料袋”；一份由院（系、所）保存并归入院（系、所）研究生教学档案。
    \item 学位类别为：工程硕士；公共管理硕士；法律硕士（非法学）；工商管理硕士；建筑学硕士；风景园林硕士；临床医学硕士；公共卫生硕士；旅游管理硕士；会计硕士；国际商务硕士；资产评估硕士；工程管理硕士；艺术硕士；工程博士；医学博士等。
    \item 本表下载区：\texttt{http://seugs.seu.edu.cn/3676/list.htm}。本表电子文档打印时用 A4 纸张，格式不变，内容较多可以加页。
\end{enumerate}

\newpage
\section*{一、学位论文开题报告}

\noindent
\begin{tabularx}{\textwidth}{|C{1.7cm}|Y|Y|Y|Y|C{1.3cm}|Y|C{1.7cm}|C{1.3cm}|Y|}
    \hline
    论文题目 & \multicolumn{9}{|l|}{\hspace{0.6em}xxx} \\ \hline
    研究方向 & \multicolumn{9}{|l|}{\hspace{0.6em}xxx} \\ \hline
    
    \multirow{2}{*}{题目来源}
    & 国家 & 部委 & 省 & 市 & 厂、矿 & 自选 & 有无合同 & 经费数 & 备注 \\ \cline{2-10}
    & & & & & & & & & \\ \hline
    
    题目类型
    & \zihao{5}理论研究 & \zihao{5}应用研究 & \zihao{5}工程技术 & \zihao{-5}跨学科研究 & \zihao{5}其他  & \multicolumn{4}{|l|}{\hspace{0.6em}xxx} \\ \hline
\end{tabularx}

\vspace{-\arrayrulewidth}
\noindent
\begin{tcolorbox}[reportbox]
    \begin{center}
    \vspace{0.45em}
    \textbf{开题报告内容}
    （具体要求见《东南大学研究生论文选题和开题报告的原则和要求》）
    \vspace{0.45em}
    \end{center}

    \vspace{-0.2em}
    \hrule
    \vspace{0.8em}

    \begingroup
        \small
        \setlength{\leftskip}{1.2em}
        \setlength{\rightskip}{1.2em}
        \setlength{\parindent}{2em}
        \setlength{\parskip}{0.45em}
        
        \noindent\textbf{一、立题依据及价值}
        
        \lipsum[1]
        
        \noindent\textbf{二、国内外研究现状}
        
        \lipsum[1-3]
        
        \noindent\textbf{三、研究内容及方法}
        
        \lipsum[1-3]
        
        \noindent\textbf{四、可行性分析}
        
        \lipsum[1-3]
        
        \noindent\textbf{五、预期成果}
        
        \lipsum[1]
        \cite{zhang2023ev}
        
        \vspace{0.8em}
        \noindent \textbf{参考文献}
        
        \vspace{0.3em}
        \noindent
        \begin{minipage}{0.90\linewidth}
        % \footnotesize
        \bibliographystyle{IEEEtran}
        \bibliography{ref}
        \end{minipage}
        
        
        \vfill
        \begin{flushright}
            \hspace*{\fill}
            \begin{minipage}{6.2cm}
                研究生签名：\emptyline[3.5cm]\\[0.8em]
                \hspace*{3.6cm}年\hspace{1.5em}月\hspace{1.5em}日
            \end{minipage}
        \end{flushright}
        \vspace{0.4em}
    \endgroup
\end{tcolorbox}

\vspace{0.5em}
（注：本页不够可附页）

\newpage
\section*{二、学位论文工作实施计划}


（一）论文的理论分析与硬件要求及其预期达到的水平与结果
\vspace{-1cm}
\begin{center}
    \noindent
    \begin{tcolorbox}[
        enhanced,
        width=0.92\textwidth,
        height=7cm,
        colback=white,
        colframe=black,
        boxrule=0.4pt,
        arc=0pt,
        outer arc=0pt,
        boxsep=0pt,
        left=10pt,
        right=10pt,
        top=10pt,
        bottom=10pt
    ]   
    \lipsum[1]
    \end{tcolorbox}
\end{center}

\vspace{0.6cm}
（二）论文工作进度与安排
\begin{center}
    \begin{tabularx}{0.92\textwidth}{|C{2.6cm}|X|C{2.5cm}|}
        \hline
        % \multicolumn{3}{|c|}{（二）论文工作进度与安排} \\ \hline
        起讫日期 & \multicolumn{1}{c|}{工作内容和要求} & 备注 \\ \hline
        202x.0x--202x.1x & s sqwerfghjmedcfvhbgvcbujsnd km,yn svcdacva & 完成 \\ \hline
        202x.0x--202x.0x & s sqwerfghjmedcfvhbgvcbujsnd km,yn svcdacva & 基本完成 \\ \hline
        202x.0x--202x.1x & s sqwerfghjmedcfvhbgvcbujsnd km,yn svcdacva & 已经初步开展 \\ \hline
        202x.0x--202x.0x & s sqwerfghjmedcfvhbgvcbujsnd km,yn svcdacva & 待完成 \\ \hline
        202x.0x--202x.1x & s sqwerfghjmedcfvhbgvcbujsnd km,yn svcdacva & 待完成 \\ \hline
        202x.0x--202x.0x & s sqwerfghjmedcfvhbgvcbujsnd km,yn svcdacva & 待完成 \\ \hline
        202x.0x--202x.1x & s sqwerfghjmedcfvhbgvcbujsnd km,yn svcdacva & 待完成 \\ \hline
    \end{tabularx}
\end{center}

\newpage
\section*{三、指导教师及审议意见}

\noindent
\begin{tabularx}{\textwidth}{|C{3.1cm}|X|}
    \hline
    
    \vlabel{8.7cm}{\makecell[c]{学校指导教师\\对开题报告的\\综合意见}}
    &
    \begin{minipage}[c][8.7cm][t]{\linewidth}
        \vspace{1em}
        
        \lipsum[1]
        
        \vfill
        \begin{flushright}
            \hspace*{\fill}
            \begin{minipage}{6.8cm}
                指导教师（签名）：\emptyline[3.5cm]\\[0.8em]
                \hspace*{4cm}年\hspace{1.5em}月\hspace{1.5em}日
            \end{minipage}
        \end{flushright}
        \vspace{0.4em}
    \end{minipage}
    \\ \hline

    \vlabel{8.7cm}{\makecell[c]{校外指导教师\\对开题报告的\\综合意见}}
    &
    \begin{minipage}[c][8.7cm][t]{\linewidth}
        \vspace{1em}
        
        \lipsum[1]
        
        \vfill
        \begin{flushright}
            \hspace*{\fill}
            \begin{minipage}{6.8cm}
                指导教师（签名）：\emptyline[3.5cm]\\[0.8em]
                \hspace*{4cm}年\hspace{1.5em}月\hspace{1.5em}日
            \end{minipage}
        \end{flushright}
        \vspace{0.4em}
    \end{minipage}
    \\ \hline

\end{tabularx}

\newpage
\vspace*{-1.0cm}

\begingroup
    % \footnotesize
    \renewcommand{\arraystretch}{0.95}
    
    \noindent
    \begin{tabularx}{\textwidth}{|C{2.6cm}|X|}
        \hline
        
        \vlabel{12.6cm}{\makecell[c]{开\\题\\报\\告\\审\\议\\情\\况\\记\\录}}
        &
        \begin{minipage}[c][12.6cm][t]{0.96\linewidth}
            \vspace{1em}
            
            1、审议小组成员（硕士至少 5 人，博士 5--7 人，其中 1 人须为校外导师）：\\[0.25em]
            
            \hspace{1.5em}组长：\emptyline[4cm]\\[0.25em]
            
            \hspace{1.5em}成员：\emptyline[4cm]\\[0.25em]
            
            2、审议小组意见\\[1.6em]
            
            3、投票表决结果\\[0.25em]
            
            \hspace{1.5em}审议小组出席 \emptyline[1.4cm] 人；通过 \emptyline[1.4cm] 人；不通过 \emptyline[1.4cm] 人。\\[0.25em]
            
            \hspace{1em}开题报告质量 \emptyline[4cm] （优、良、中、通过）\\[0.25em]
            
            4、审议小组组长（签名）：\emptyline[4cm]\\[0.25em]
            
            \hspace{1.5em}审议小组成员（签名）：\emptyline[4cm]\\[0.25em]
            
            \vfill
            \begin{flushright}
                年\hspace{1.5em}月\hspace{1.5em}日
            \end{flushright}
            
            \vspace{0.1em}
        \end{minipage}
        \\ \hline
        
        \vlabel{3.0cm}{院（系、所）意见}
        &
        \begin{minipage}[c][2.0cm][t]{0.96\linewidth}
            \vfill
            \hspace{1.5em} \\
            \begin{flushright}
                院（系、所）负责人签名（或印章）：\emptyline[4.0cm]\\[0.35em]
                年\hspace{1.5em}月\hspace{1.5em}日
            \end{flushright}
            \vspace{0.1em}
        \end{minipage}
        \\ \hline
        
        \vlabel{1.25cm}{备注：}
        &
        \begin{minipage}[c][2cm][t]{0.96\linewidth}
            \hspace{1cm}
        \end{minipage}
        \\ \hline
        
    \end{tabularx}

\endgroup

\end{document}